% ==============================================================================
% CONCLUSIONES Y RECOMENDACIONES
% ==============================================================================
\newpage
\section*{\normalfont\textbf{Conclusiones y Recomendaciones}}
\addcontentsline{toc}{section}{Conclusiones y Recomendaciones}

\subsection*{\normalfont\textbf{Conclusiones}}
\addcontentsline{toc}{subsection}{Conclusiones}

% [TEXTO TEMPORAL]: Redactar las conclusiones definitivas una vez finalizada la implementación y las pruebas del sistema.

A partir del perfil de proyecto desarrollado, se concluye que el proceso actual de postulación y adjudicación de la Beca Comedor en la Dirección Universitaria de Bienestar Social y Salud (DUBSS) de la Universidad Autónoma Gabriel René Moreno presenta una ineficiencia estructural significativa, derivada de la revisión manual de un volumen de solicitudes que supera ampliamente la capacidad operativa del personal administrativo. El diseño del sistema propuesto, sustentado en el Proceso Unificado de Desarrollo de Software y en un modelo de precalificación basado en Inteligencia Artificial, plantea una solución viable para automatizar la evaluación socioeconómica y garantizar una asignación más objetiva y transparente de los cupos disponibles. El marco teórico, la especificación de requisitos según IEEE 830 y el modelado UML desarrollados hasta este punto constituyen una base sólida para las siguientes fases de diseño detallado, implementación y pruebas del sistema.

\subsection*{\normalfont\textbf{Recomendaciones}}
\addcontentsline{toc}{subsection}{Recomendaciones}

% [TEXTO TEMPORAL]: Redactar las recomendaciones definitivas una vez finalizada la implementación y las pruebas del sistema.

Se recomienda continuar con las fases de Elaboración y Construcción del PUDS conforme al cronograma establecido, dando prioridad a la validación del modelo de precalificación con datos históricos reales de la DUBSS antes de su despliegue en producción. Asimismo, se sugiere planificar con anticipación la capacitación del personal administrativo y de auditoría en el uso del nuevo sistema, y establecer desde etapas tempranas los protocolos de seguridad y resguardo de la información socioeconómica de los postulantes, dada la sensibilidad de dichos datos.
